\chapter{Fonctions : définition, graphe et opérations}
\label{workbook:chapter:03}
\section{La notion de fonction}
\label{workbook:section:03.01}
\begin{exerciseblock}{La notion de fonction}
\item Parmi les relations suivantes, lesquelles définissent une fonction de $x$ vers $y$ : $y=2x-1$; $x=y^2$; $y=\pm\sqrt{x}$; $y=|x|$? Justifier.
\item Soit $f(x)=\dfrac{\sqrt{x+2}}{x-1}$. Déterminer précisément $\dom(f)$.
\item Une machine transforme une température $C$ en $F=1{,}8C+32$. Identifier entrée, sortie, variable indépendante et unité de chacune.
\end{exerciseblock}
\section{Graphe et lecture graphique}
\label{workbook:section:03.02}
\begin{exerciseblock}{Graphe et lecture graphique}
\item Pour \(f(x)=x^2-4x+3=(x-1)(x-3)=(x-2)^2-1\), utiliser les formes données pour déterminer les intersections avec les axes, le sommet et l'axe de symétrie, puis esquisser le graphe.
\item Décrire comment lire graphiquement les solutions de $f(x)=2$ et de $f(x)>2$.
\item Un graphe de température est défini pour $0\le t\le24$. Expliquer la différence entre domaine mathématique naturel et domaine imposé par le contexte.
\end{exerciseblock}
\section{Transformations de graphes}
\label{workbook:section:03.03}
\begin{exerciseblock}{Transformations de graphes}
\item À partir de $y=x^2$, décrire puis esquisser $y=-2(x-3)^2+1$.
\item Si $(a,b)$ est sur le graphe de $f$, donner les points correspondants sur les graphes de $f(x)+4$, $f(x-2)$, $-f(x)$ et $f(-x)$.
\item Comparer les effets de $f(2x)$ et de $2f(x)$ sur le graphe, à l'aide de \(f(x)=|x-1|\). Vérifier en particulier l'abscisse du zéro de chaque graphe.
\end{exerciseblock}
\section{Opérations et composition}
\label{workbook:section:03.04}
\begin{exerciseblock}{Opérations et composition}
\item Pour $f(x)=\sqrt{x+1}$ et $g(x)=1/(x-2)$, déterminer les domaines de $f+g$, $fg$ et $f/g$.
\item Avec $f(x)=2x-3$ et $g(x)=x^2+1$, calculer $f\circ g$ et $g\circ f$, puis montrer qu'elles diffèrent.
\item Une taxe de $15\%$ est appliquée puis un rabais de $20\%$. Modéliser les deux opérations comme fonctions et comparer les deux ordres.
\end{exerciseblock}
\section{Fonction réciproque}
\label{workbook:section:03.05}
\begin{exerciseblock}{Fonction réciproque}
\item Trouver la réciproque de $f(x)=3x-7$ et vérifier les deux compositions.
\item Expliquer pourquoi $x^2$ n'a pas de réciproque sur $\R$, puis en construire une après restriction à $[0,+\infty[$.
\item Déterminer la réciproque de $f(x)=\dfrac{2x+1}{x-3}$ en précisant domaine et image.
\end{exerciseblock}
\section{Situation, formule, tableau et graphe}
\label{workbook:section:03.06}
\begin{exerciseblock}{Situation, formule, tableau et graphe}
\item Un réservoir contient $120$ L et perd $8$ L par minute. Construire la formule, un tableau pour $t=0,5,10,15$, et décrire le graphe.
\item À partir du tableau $(0,4)$, $(1,7)$, $(2,10)$, $(3,13)$, proposer une formule et justifier le modèle.
\item Inventer un contexte correspondant à \(f(t)=50-3t\), pour \(0\le t\le50/3\), et interpréter \(50\), \(3\) et \(t\) avec leurs unités.
\end{exerciseblock}
\section{Lire une fonction à partir de son graphe}
\label{workbook:section:03.07}
\begin{exerciseblock}{Lire une fonction à partir de son graphe}
\item Dessiner, sans lever le crayon, le graphe d'une fonction sur \([-3,4]\) ayant deux zéros, un maximum local et un minimum global, puis indiquer ces caractéristiques.
\item Sur un graphe donné, expliquer comment distinguer maximum local, maximum global, zéro et ordonnée à l'origine.
\item Une courbe coupe la droite $y=x$ en trois points. Que signifient ces points pour l'équation $f(x)=x$?
\end{exerciseblock}
\section{Analyser une fonction sous plusieurs représentations}
\label{workbook:section:03.08}
\begin{exerciseblock}{Analyser une fonction sous plusieurs représentations}
\item Pour $f(x)=|x-2|+1$, produire formule par morceaux, tableau de cinq valeurs et description graphique.
\item Une fonction est donnée par $f(-2)=1$, $f(0)=3$, $f(2)=1$. Proposer deux modèles différents compatibles et expliquer pourquoi trois points ne déterminent pas toujours une fonction unique.
\item Pour \(f(x)=x^2-4x+3=(x-1)(x-3)=(x-2)^2-1\), comparer les informations visibles dans les trois écritures données : valeur en zéro, zéros et sommet.
\end{exerciseblock}
\section*{Synthèse du chapitre}
\begin{synthesisbox}
\item Écrire \(f(x)=|x-1|\) par morceaux. Tracer ses deux demi-droites et expliquer pourquoi elles se rejoignent en \(x=1\) tout en formant un coin.
\item Deux fonctions $f$ et $g$ ont le même graphe visible sur $[-2,2]$. Expliquer pourquoi cela ne suffit pas pour conclure qu'elles sont égales comme fonctions.
\item Modéliser une conversion composée : dollars canadiens vers euros, puis ajout de frais fixes. Identifier les unités à chaque étape et écrire la composition.
\end{synthesisbox}
